\documentclass[border=10pt]{standalone}
\usepackage[no-math]{fontspec}
\usepackage[brazil]{babel}
\usepackage{amsmath,amssymb}
\IfFileExists{C:/Windows/Fonts/Lato_Regular.ttf}{%
  \setmainfont{Lato}[
    Path = C:/Windows/Fonts/,
    UprightFont = Lato_Regular.ttf,
    BoldFont = Lato_Bold.ttf,
    ItalicFont = Lato_Italic.ttf,
    BoldItalicFont = Lato_BoldItalic.ttf
  ]
  \setsansfont{Lato}[
    Path = C:/Windows/Fonts/,
    UprightFont = Lato_Regular.ttf,
    BoldFont = Lato_Bold.ttf,
    ItalicFont = Lato_Italic.ttf,
    BoldItalicFont = Lato_BoldItalic.ttf
  ]
}{%
  \RequirePackage[default]{lato}
}
\renewcommand{\familydefault}{\sfdefault}

\usepackage{xcolor}
\usepackage{tikz}
\usepackage{pgfplots}
\pgfplotsset{compat=1.18}
\usetikzlibrary{arrows.meta,positioning,shapes.geometric,fit,backgrounds,calc,decorations.pathreplacing,patterns}

% Paleta institucional UEFS / Bioengenharia de Solos
\definecolor{bsNavy}{HTML}{1F4268}
\definecolor{bsNavyDark}{HTML}{152E4A}
\definecolor{bsGold}{HTML}{B38A3D}
\definecolor{bsGoldDark}{HTML}{8C6927}
\definecolor{bsTeal}{HTML}{0E7C7B}
\definecolor{bsForest}{HTML}{2D6A4F}
\definecolor{bsAlert}{HTML}{9E2A2B}
\definecolor{bsEarth}{HTML}{8C5A3C}
\definecolor{bsWater}{HTML}{2E86AB}
\definecolor{bsInk}{HTML}{26313D}
\definecolor{bsMuted}{HTML}{5B6573}

\definecolor{bgNavy}{HTML}{EDF2F7}
\definecolor{bgGold}{HTML}{F8F3E9}
\definecolor{bgTeal}{HTML}{EBF6F5}
\definecolor{bgForest}{HTML}{EEF5F0}
\definecolor{bgAlert}{HTML}{FDF2F2}

\tikzset{
  >=Stealth,
  bsBox/.style={
    draw=#1, fill=white, rounded corners=5pt, line width=1.15pt,
    inner sep=7pt, align=center, font=\footnotesize\sffamily, text=bsInk
  },
  bsTint/.style={
    draw=#1, fill=#1!8!white, rounded corners=5pt, line width=1.15pt,
    inner sep=7pt, align=center, font=\footnotesize\sffamily, text=bsInk
  },
  bsBadge/.style={
    fill=#1, text=white, rounded corners=9pt, inner xsep=8pt, inner ysep=3pt,
    font=\scriptsize\bfseries\sffamily
  },
  bsTitle/.style={font=\normalsize\bfseries\sffamily, text=bsNavy, align=center},
  bsSub/.style={font=\footnotesize\sffamily, text=bsMuted, align=center},
  bsArrow/.style={->, >=Stealth, line width=1.35pt},
  bsDArrow/.style={<->, >=Stealth, line width=1.35pt}
}


\begin{document}
\begin{tikzpicture}
  \hyphenpenalty=10000\exhyphenpenalty=10000

  % Título
  \node[bsTitle] at (1.5, 5.50)
    {Árvore de Decisão para Seleção de Técnicas no Projeto Integrador};
  \node[bsSub] at (1.5, 5.02)
    {Do diagnóstico de sítio ao pacote tecnológico compatível com o processo erosivo dominante};

  % Raiz esquerda
  \node[bsTint=bsNavy, text width=3.65cm, minimum height=2.65cm] (root) at (-5.65, 0.85) {
    \vspace{0.25cm}
    \textbf{\color{bsNavy}Diagnóstico do Sítio}\\[5pt]
    \scriptsize
    \textbullet\ Declividade ($S$) e comprimento\\
    \textbullet\ Textura e coesão do solo\\
    \textbullet\ Vazão de projeto ($Q$)\\
    \textbullet\ Mecanismo de falha dominante
  };
  \node[bsBadge=bsNavy] at (root.north) {PONTO DE PARTIDA};

  % 4 Regimes Diagnósticos (Coluna Central, espaçamento vertical generoso)
  \node[bsTint=bsGold, text width=4.25cm, minimum height=1.35cm] (d1) at (-0.15, 3.65) {
    \textbf{\color{bsGoldDark}1. Erosão Laminar / Sulco Raso}\\
    \scriptsize Escoamento difuso em solo exposto
  };
  \node[bsTint=bsAlert, text width=4.25cm, minimum height=1.35cm] (d2) at (-0.15, 1.80) {
    \textbf{\color{bsAlert}2. Ravina / Voçoroca Ativa}\\
    \scriptsize Incisão concentrada e avanço remontante
  };
  \node[bsTint=bsEarth, text width=4.25cm, minimum height=1.35cm] (d3) at (-0.15, -0.05) {
    \textbf{\color{bsEarth}3. Talude Íngreme / Aterro Instável}\\
    \scriptsize Risco de ruptura por poropressão/peso
  };
  \node[bsTint=bsWater, text width=4.25cm, minimum height=1.35cm] (d4) at (-0.15, -1.90) {
    \textbf{\color{bsWater}4. Margem Fluvial em Solapamento}\\
    \scriptsize Tensão cisalhante $\tau_0 > \tau_{\text{adm}}$ na base
  };

  % 4 Pacotes de Solução (Coluna Direita)
  \node[bsBox=bsGold, text width=5.65cm, minimum height=1.35cm] (s1) at (6.05, 3.65) {
    \textbf{\color{bsGoldDark}Hidrossemeadura + Biomantas + Cordões}\\
    \scriptsize Cobertura imediata contra impacto da gota e redução de velocidade superficial
  };
  \node[bsBox=bsAlert, text width=5.65cm, minimum height=1.35cm] (s2) at (6.05, 1.80) {
    \textbf{\color{bsAlert}Drenagem Montante + Barraginhas + Paliçadas em Série}\\
    \scriptsize Desvio do fluxo, dissipação em degraus e retenção de sedimentos
  };
  \node[bsBox=bsEarth, text width=5.65cm, minimum height=1.35cm] (s3) at (6.05, -0.05) {
    \textbf{\color{bsEarth}Dreno Espinha-de-Peixe + Fascinas + Parede Krainer}\\
    \scriptsize Rebaixamento de poropressão e contenção mista com enraizamento
  };
  \node[bsBox=bsWater, text width=5.65cm, minimum height=1.35cm] (s4) at (6.05, -1.90) {
    \textbf{\color{bsWater}Enrocamento Vegetado + Gabião Vivo + Mata Ciliar}\\
    \scriptsize Proteção basal rígida-flexível e restauração da faixa ciliar
  };

  % Conectores tronco -> regimes
  \coordinate (split) at (-3.15, 0.85);
  \draw[bsNavy, line width=1.35pt] (root.east) -- (split);
  \draw[bsArrow, bsGoldDark] (split) |- (d1.west);
  \draw[bsArrow, bsAlert]    (split) |- (d2.west);
  \draw[bsArrow, bsEarth]    (split) |- (d3.west);
  \draw[bsArrow, bsWater]    (split) |- (d4.west);

  % Conectores regimes -> soluções
  \draw[bsArrow, bsGoldDark] (d1.east) -- (s1.west);
  \draw[bsArrow, bsAlert]    (d2.east) -- (s2.west);
  \draw[bsArrow, bsEarth]    (d3.east) -- (s3.west);
  \draw[bsArrow, bsWater]    (d4.east) -- (s4.west);

\end{tikzpicture}
\end{document}
